\documentclass[10pt,a4paper]{amsart}
\usepackage[margin=19mm]{geometry}
\usepackage{amsmath,amssymb,mathtools}
\usepackage[scr=rsfs,cal=euler]{mathalfa}
\usepackage{xfrac}
\usepackage[unicode,hidelinks,bookmarks=false]{hyperref}
\newcommand{\ie}{i.e.\ }
\newcommand{\cf}{cf.\ }
\newcommand{\resp}{respectively\ }
\newcommand{\Subsection}{\subsection}
\providecommand{\nobreakdash}{\nobreak\hbox{-}}
\setlength{\emergencystretch}{2em}
\multlinegap0pt
\usepackage{mdframed}
\usepackage{mdframed}
\usepackage{mdframed}
\usepackage{mdframed}
\usepackage{amssymb}
\usepackage{mathtools}
\usepackage{bm}
\usepackage{mdframed}
\usepackage{xcolor}
\usepackage{xfrac}
\usepackage{dsfont}
\newtheorem{proposition}{Proposition}
\newtheorem{lemma}{Lemma}
\newcommand{\Ent}{\textup{\textsf{Ent}}}
\newcommand{\FOM}{\textsf{FOM}}
\newcommand{\Fh}{\Ent_d}
\newcommand{\Mtd}{\mathscr{M}}
\newcommand{\Orc}{\mathscr{O}}
\newcommand{\R}{\mathbb{R}}
\newcommand{\Risk}{\mathsf{Risk}}
\newcommand{\Simp}{\Delta_d}
\newcommand{\bnabla}{\nabla}
\newcommand{\ones}{\mathsf{1}_d}
\newcommand{\ri}{\operatorname{ri}}
\renewcommand{\top}{\mathsf{T}}
\newcommand{\wt}{\widetilde}
\title{Entropy-Smooth Convex Optimization Cannot Be Accelerated}
\author{Jacob M. Aguirre, Dmitrii M. Ostrovskii}
\date{Source: arXiv:2607.27476v1 (CC BY 4.0)}
\begin{document}
\maketitle

\noindent\emph{Source and licence.} Entropy-Smooth Convex Optimization Cannot Be Accelerated. Version arXiv:2607.27476v1 (CC BY 4.0), distributed by arXiv under the Creative Commons Attribution 4.0 International licence, http://creativecommons.org/licenses/by/4.0/, as recorded in the arXiv licence metadata for this exact version and shown on the abstract page https://arxiv.org/abs/2607.27476v1. Full licence notice: \texttt{licenses/arxiv-2607.27476-CC-BY-4.0.txt}.\par\smallskip

\noindent\textit{Editorial scope.} Counted unit: the article's proposition prop:emulation (source0.tex lines 1847--1853). The article's complete original proof of it is delivered with it (source0.tex lines 1855--1931, one \texttt{proof} environment of 77 lines). No mathematics is written, repaired or re-derived: the statement and the proof are the article's own lines, in the article's own notation, and no retained line is substituted or re-flowed.  Retained uncounted as the source-local context the proof needs: lines 535--538; lines 793--802; lines 1152--1155; lines 1309--1312; lines 1719--1723; lines 1727--1730; lines 1788--1791; lines 1812--1818; lines 1823--1832.  Nothing was omitted from any retained span, so the package carries no omission record.  No retained line contains a citation, so the wrapper carries no bibliography and no bibliography entry.  No retained line contains a figure environment, an image-inclusion command or drawing code, so no image is delivered and the package declares no asset dependency.  Editorial formatting only, in addition: an AMS \texttt{amsart} preamble that restates the article's own declarations and macros used by the retained lines, namely the environments \texttt{proposition}, \texttt{lemma} on separate counters, together with the standard packages those lines need.  The retained environments print in this document as Proposition 1, Lemma 1 in order of appearance, whereas the article prints the same items with the numbering of its own full document.  The complete native source of the article as distributed in the arXiv e-print for this exact version is bundled unchanged as \texttt{sources/206350/source0.tex}.
\begin{equation}
\label{eq:grad-tangent}
\bnabla f(x) := \nabla_\star f(x) - d^{-1}  \langle \nabla_\star f(x), \ones\rangle \ones, \qquad x \in \ri(\Simp).
\end{equation}

\begin{equation}
\label{eq:minimax-risk}
\Risk_{d}(T,L) \; 
:=
\inf_{\Mtd^{\vphantom{2^L}} \,\in\, \FOM_{d}(T)} 
\quad
\sup_{f \,\in\, \Fh(L)} 
\quad
\left\{ f(x_{T+1}^\Mtd(f))- \min_{x \in \Simp} f(x) \right\}.
\end{equation}

\begin{equation}
\label{eq:oracle-answers}
\Orc_{f_t}(x_t) = (f_t(x_t), \nabla f_t(x_t)),
\end{equation}

\begin{proposition}[Transcript consistency]
\label{prop:transcript}
The queries~$x_1,\dots,x_T$, the oracle answers returned in~\eqref{eq:oracle-answers}, and the reported point~$x_{T+1}$ coincide with the transcript of~$\Mtd$ run on the objective~$f_{T+1}$.
\end{proposition}

\begin{equation}
\label{eq:conic-extension}
f^+(x) = s(x) f (s(x)^{-1}x).
%+ \langle \ones, x \rangle - 1.
\end{equation}

\begin{equation}
\label{eq:grad-conic}
\nabla_\star f^+(x) = \bnabla f(x) + (f(x) - \langle \bnabla f(x), x \rangle) \ones.
\end{equation}

\begin{lemma}
\label{lem:conic-extension}
The 1-homogeneous extension $f^+(x)$ of any~$f \in \Fh(L)$, cf.~\eqref{eq:conic-extension}, belongs to~$\Fh^+(L)$.
\end{lemma}

\begin{equation}
\label{eq:method-maps-conic}
\begin{aligned}
\Psi_t: \R^{d(t-1)}  &\to \ri(\Simp), \qquad t \in [T];\\
\bar  \Psi_{T+1}: \R^{dT} &\to \Simp,
\end{aligned}
\end{equation}

\begin{equation}
\label{eq:minimax-risk-conic}
\Risk_{d}^+(T,L) \;\; 
:=
\inf_{\tilde\Mtd^{\vphantom{{{2^2}^2}}} \,\in\, \FOM_{d}^+(T)} 
\quad
\sup_{\tilde f \,\in\, \Fh^+(L)} 
\quad
\left\{ \tilde f(x_{T+1}^{\tilde \Mtd}(\tilde f))- \min_{x \in \Simp} \tilde f(x) \right\},
\end{equation}

\begin{proposition}
\label{prop:emulation}
The minimax risks defined in~\eqref{eq:minimax-risk-conic} and~\eqref{eq:minimax-risk} satisfy
$
\Risk_{d}^+(T,L) \ge \Risk_{d}(T,L).
$
\end{proposition}

\begin{proof}
Replacing~$\Ent_d^+(L)$ with the smaller (due to Lemma~\ref{lem:conic-extension}) class~$\{f^+: f \in \Fh(L)\}$, we get
\[
\begin{aligned}
\Risk_{d}^+(T,L) 
&\ge 
\inf_{\tilde\Mtd^{\vphantom{{{2^2}^2}}} \,\in\, \FOM_{d}^+(T)} 
\quad
\sup_{f \,\in\, \Fh(L) \vphantom{\wt \Fh(L)}} 
\quad
\left\{  f^+(x_{T+1}^{\tilde \Mtd}(f^+))- \min_{x \in \Simp} f^+(x) \right\} \\ 
&= 
\inf_{\tilde\Mtd^{\vphantom{{{2^2}^2}}} \,\in\, \FOM_{d}^+(T)} 
\quad
\sup_{f \,\in\, \Fh(L) \vphantom{\wt \Fh(L)}} 
\quad
\left\{  f(x_{T+1}^{\tilde \Mtd}(f^+))- \min_{x \in \Simp} f(x) \right\}
\ge \Risk_{d}(T,L).
\end{aligned}
\]
%Here,~$f^+$ denotes the conic extension of~$f$, and 
Here the identity holds since the mappings in~\eqref{eq:method-maps-conic} output points in~$\Simp$ (where~$f^+ = f$). 
For the last inequality, we use~\eqref{eq:grad-conic} and observe
%\[
%\nabla f_+(x) 
%= \bnabla f(x) + (f(x) - \langle \bnabla f(x), x - \tfrac{1}{d} \ones \rangle) \ones.
%\]
%~\eqref{eq:grad-conic}--\eqref{eq:grad-tangent} to express~$\nabla f^+(x)$ in terms of~$\Orc_f(x) = (f(x),\bnabla f(x))$:
%\begin{equation*}
%\nabla f^+(x) = \bnabla f(x) + (f(x) - \bnabla f(x)^\top x) \ones.
%\end{equation*}
that the execution on~$f^+$ of~$\smash{\tilde\Mtd \in \FOM_d^+(T)}$, as per~\eqref{eq:method-maps-conic},  produces the same query sequence
\iffalse
\[
\begin{aligned}
\tilde x_1 &= \Psi_1, \\
\tilde x_2 &= 
\Psi_2 \left(\bnabla f(\tilde x_1) + \big(f(\tilde x_1) -\langle \bnabla f(\tilde x_1),  \tilde x_1 - d^{-1} \ones \rangle \big) \ones \right), \\
%\;=\; \Phi_2^{\vphantom{\Mtd}}(\Orc_f(x_1^\Mtd)), \\
\tilde x_{T+1} &= \bar\Psi_{T+1}
\Big( \bnabla f(\tilde x_1) + \big(f(\tilde x_1) \;- \langle \bnabla f(\tilde x_1),  \tilde x_1 - d^{-1} \ones \rangle \big) \ones, \dots, \\
&\qquad \qquad \, \bnabla f(\tilde x_T) + \big(f(\tilde x_T) -\langle \bnabla f(\tilde x_T),  \tilde x_T - d^{-1} \ones \rangle \big) \ones \Big),
\end{aligned}
\]
\fi
as the execution on~$f$ of the method
$
\Mtd \in \FOM_d(T),
$
defined by
%whose mappings~$\Phi_t$ are inductively defined  by 
\newcommand{\bg}{\xi}
\[
\begin{aligned}
\Phi_1 &:= \Psi_1, \\
\Phi_2(v_1, \bg_1) 
&:= \Psi_2 \left(\bg_1 + \big(v_1 -\langle \bg_1,  \Phi_1 \rangle \big) \ones \right), \\
\Phi_3(v_1,\bg_1, v_2, \bg_2) 
&:= \Psi_3 
		\left(
			\bg_1 + \big(v_1 -\langle \bg_1,  \Phi_1 \rangle \big) \ones,  \;\;
			\bg_2 + \big(v_2 -\langle \bg_2,  \Phi_2(v_1,\bg_1) \rangle \big) \ones
		\right), \\
&\;\;\vdots \vspace{-0.1cm} \\ \vspace{-0.1cm}
\bar\Phi_{T+1}(v_1,\bg_1, \dots, v_T, \bg_T) \\
:= \bar\Psi_{T+1} 
		\big(
		\bg_1 + \big(v_1 - &\langle \bg_1, \Phi_1 \rangle \big) \ones, \;\dots,\;
		\bg_T + \big(v_T -\langle \bg_T,  \Phi_T(v_1,\bg_1, \dots, v_{T-1},\bg_{T-1}) \rangle \big) \ones
		\big).
\end{aligned}
\]
%whose~$\Phi_t$ is defined as the composition of~$\Psi_t$ and 
%depends on~$\nabla f$ only through~$\bnabla f$.)
This can be verified by induction over~$t$, in the same way as in the proof of Proposition~\ref{prop:transcript}.
This fact implies the desired inequality, since in~$\Risk_{d}(T,L)$ the infimum is over {\em all} methods in~$\FOM_d(T)$. 
\end{proof}

\end{document}
